\documentclass[11pt,a4paper]{article}

% ====== Core Packages ======
\usepackage[margin=1in]{geometry}    % Set page margins
\usepackage{amsmath, amsthm, amssymb} % Math formulas and proof environments
\usepackage{graphicx}                % Insert images (experimental result charts)
\usepackage{algorithm}               % Pseudocode float
\usepackage{algpseudocode}           % Pseudocode formatting
\usepackage{hyperref}                % Hyperlinks and bookmarks
\usepackage{listings}                % Insert code snippets
\usepackage{xcolor}                  % Color support
\usepackage{booktabs}                % Three-line table (used for cost analysis tables)

% ====== Theorem and Proof Environments Setup ======
\newtheorem{theorem}{Theorem}[section]
\newtheorem{lemma}[theorem]{Lemma}
\newtheorem{invariant}{Loop Invariant}[section] % Custom environment defined specifically for loop invariants

% ====== Code Highlighting Setup (Python) ======
\lstset{
    language=Python,
    basicstyle=\ttfamily\small,
    keywordstyle=\color{blue}\bfseries,
    stringstyle=\color{red},
    commentstyle=\color{green!50!black},
    numbers=left,
    numberstyle=\tiny\color{gray},
    stepnumber=1,
    frame=single,
    breaklines=true
}

% ====== Title and Author Information ======
\title{\textbf{Red-Black Tree: Design, Implementation, and Analysis} \\ 
\large CSCI 4041 Midterm Project Report}
\author{
  Team Member 1 \and 
  Team Member 2 \and 
  Team Member 3 \and 
  Team Member 4
}
\date{\today}

\begin{document}

\maketitle

\begin{abstract}
This report presents a comprehensive study of the Red-Black Tree, a self-balancing binary search tree. We detail the theoretical design, pseudocode, and line-by-line cost analysis of insertion and deletion operations. Furthermore, we provide rigorous mathematical proofs for the correctness of these algorithms using loop invariants. Finally, we discuss our Python implementation and present empirical performance data confirming the $O(\log n)$ asymptotic bounds.
\end{abstract}

% ==========================================
\section{Introduction}
% ==========================================
A red-black tree is a binary search tree with one extra bit of storage per node: its color, which can be either RED or BLACK. By constraining the node colors on any simple path from the root to a leaf, red-black trees ensure that no such path is more than twice as long as any other, so that the tree is approximately balanced. Indeed, as we're about to see, the height of a red-black tree with $n$ keys is at most $2\log{(n+1)}$, which is $O\log{(n)}$. \newline
\indent Each node of the tree now contains the attributes \textit{color}, \textit{key}, \textit{left}, \textit{right}, and $p$. If a child or the parent of a node does not exist, the corresponding pointer attribute of the node contains the value NIL. Think of these NILs as pointers to leaves (external nodes) of the binary search tree and the normal, key-bearing nodes as internal nodes of the tree.\newline
\indent Overall A red-black tree is a binary search tree that satisfies the following \textbf{\textit{red-black properties}}:

\begin{enumerate}
    \item Every node is either red or black.
    \item The root is black.
    \item Every leaf (\textsc{nil}) is black.
    \item If a node is red, then both its children are black.
    \item For each node, all simple paths from the node to descendant leaves contain the same number of black nodes.
\end{enumerate}
% ==========================================
\section{Design and Paper Analysis}
% ==========================================
This section outlines the pseudocode and complexity analysis for the primary operations.

\subsection{Insertion Algorithm}
Here we present the insertion algorithm and its line-by-line cost.

% Pseudocode example
\begin{algorithm}
\caption{Red-Black Tree Insert(T, z)}\label{alg:rb_insert}
\begin{algorithmic}[1]
\State $y \gets T.nil$ \Comment{Cost: $c_1$}
\State $x \gets T.root$ \Comment{Cost: $c_2$}
\While{$x \neq T.nil$} \Comment{Executes $O(\log n)$ times}
    \State $y \gets x$
    \If{$z.key < x.key$}
        \State $x \gets x.left$
    \Else
        \State $x \gets x.right$
    \EndIf
\EndWhile
\State $z.p \gets y$
\State \dots \Comment{Continue with balancing logic}
\State \Call{RB-Insert-Fixup}{T, z} \Comment{Cost: $O(\log n)$}
\end{algorithmic}
\end{algorithm}

\subsection{Correctness Proof: Loop Invariant}
To prove the correctness of our insertion (specifically the fix-up phase), we establish the following loop invariant.

\begin{invariant}
At the start of each iteration of the while loop in \texttt{RB-Insert-Fixup}, the following hold:
\begin{enumerate}
    \item Node $z$ is red.
    \item If $z.p$ is the root, then $z.p$ is black.
    \item If there is a violation of the red-black properties, it is at most one violation, and it is either property 2 or property 4.
\end{enumerate}
\end{invariant}

\begin{proof}
\textbf{Initialization:} Prior to the first iteration... \\
\textbf{Maintenance:} Assume the invariant holds before an iteration. We consider the various cases (e.g., $z$'s uncle is red vs. black)... \\
\textbf{Termination:} When the loop terminates, $z.p$ is black. Thus, no red nodes are adjacent, and the properties are restored.
\end{proof}

\subsection{Asymptotic Runtime Analysis}
The best-case runtime for insertion is $\Omega(1)$ (e.g., empty tree). The worst-case runtime is bounded by $O(\log n)$ because the height of a Red-Black tree of $n$ nodes is at most $2\log(n+1)$.

% ==========================================
\section{Implementation Details}
% ==========================================
In this section, we discuss translating our theoretical design into Python, adhering strictly to the requirements without using generative AI tools. We utilize a \texttt{NIL} sentinel node to simplify boundary conditions.

% Insert a small key code snippet for demonstration
\begin{lstlisting}[caption={Example of Initialization of BRT and Node}, label=lst:Initialization]
class node:
    def __init__(self, key, left = None, right = None, parent = None, color = RED):
        self.key = key
        self.p = parent
        self.left = left
        self.right = right
        self.color = color 

class RedBlackTree:
    def __init__(self):
        self.NIL = node(key = None, color = BLACK)
        self.NIL.left = None
        self.NIL.right = None
        self.root = self.NIL
\end{lstlisting}

\begin{lstlisting}[caption={Example of Left Rotation in Python}, label=lst:left_rotate]
def left_rotate(self, x):
    # Initialization
    y = x.right
    x.right = y.left

    # Connect y left child to x Connect y to x parent
    if y.left != self.NIL:
        y.left.p = x
    y.p = x.p
    
    # If x is root, set y to root, else y replace x place
    if x.p == self.NIL: 
        self.root = y
    elif x.p.left == x:
        x.p.left = y
    else:
        x.p.right = y
    
    # Reconnect x and y, y is parent
    y.left = x
    x.p = y
\end{lstlisting}

\begin{lstlisting}[caption={Example of Insert},label=lst:insert]
def insert(self, key):
    # Create a new node
    z = node(key)
    # Find the place to insert z
    x = self.root
    y = self.NIL
    while x != self.NIL:
        # Save z's parent
        y = x
        if z.key < x.key:
            x = x.left
        else:
            x = x.right
    z.p = y
    # If there is no node in tree, z be the root
    if y == self.NIL: 
        self.root = z
    elif z.key < y.key:
        y.left = z
    else:
        y.right = z
    
    # Setting the z child to NIL
    z.left = self.NIL
    z.right = self.NIL

    self.insert_fixup(z)
\end{lstlisting}

\begin{lstlisting}[caption={Example of Insert Fix up},label=lst:fixup]
def insert_fixup(self, z):
    while z.p.color == RED:
        if z.p == z.p.p.left: 
            y = z.p.p.right # Check the uncle's color
            # Case 1, if uncle is RED, 
            if y.color == RED:
                z.p.color = BLACK # Let parent and uncle be BLACK
                y.color = BLACK
                # Set grandparent to RED to maintain subtree's "Black height"
                z.p.p.color = RED
                z = z.p.p
            else:
                # Case 2, z is right child
                # Left rotate z and prepare for Case 3
                if z == z.p.right:
                    z = z.p
                    self.left_rotate(z)
                # Case 3, z is left child
                # Maintain the subtree's height and Right rotate z grandparent 
                z.p.color = BLACK
                z.p.p.color = RED
                self.right_rotate(z.p.p)
        else:
            y = z.p.p.left
            if y.color == RED:
                z.p.color = BLACK
                y.color = BLACK
                z.p.p.color = RED
                z = z.p.p
            else:
                if z == z.p.left:
                    z = z.p
                    self.right_rotate(z)
                z.p.color = BLACK
                z.p.p.color = RED
                self.left_rotate(z.p.p)
    self.root.color = BLACK
\end{lstlisting}
\begin{lstlisting}[caption={Example of Transplant},label=lst:Transplant]
def transplant (self, u, v):
    # Transplant the node u with target node v
    if u.p == self.NIL:
        self.root = v
    elif u == u.p.left:
        u.p.left = v
    else:
        u.p.right = v
    v.p = u.p
\end{lstlisting}

\begin{lstlisting}[caption={Example of Delete}, label = lst:Delete]
def delete(self, z):
    y = z
    y_original_color = y.color
    # If 
    if z.left == self.NIL:
        x = z.right
        self.transplant(z, z.right) # Replace by right child
    elif z.right == self.NIL:
        x = z.left
        self.transplant(z, z.left) # Replace by left child
    else:
        y = self.min(z.right) # Find the successor node
        y_original_color = y.color
        # x is y's right child, which will replace y's original position.
        x = y.right
        if y!=z.right:
            # If y is deeper in the tree, replace y with x
            self.transplant(y, y.right)
            y.right = z.right
            y.right.p = y
        else:
            # If y is z's direct child
            # Update x's parent pointer (important if x is NIL)
            x.p = y
        # Replace z with y in the tree
        self.transplant(z, y)
        # Attach z's left subtree to y and inherit z's color
        y.left = z.left
        y.left.p = y
        y.color = z.color
    
    if y_original_color == BLACK:
        self.delete_fix(x)
\end{lstlisting}

\begin{lstlisting}[caption={Example of Delete Fix up}, label = lst:deleteFixUp]
    
\end{lstlisting}

% ==========================================
\section{Experimentation and Numerical Results}
% ==========================================
We performed timing tests for best, worst, and average cases. 

\subsection{Test Case Generation}
We generated sorted data (worst-case for naive BST, triggers maximum rebalancing in RBT) and randomized data (average case). 

\subsection{Performance Graphs}
As shown in Figure \ref{fig:runtime}, the insertion time scales logarithmically.

% Placeholder for image insertion (upload your experiment result charts to Overleaf)
\begin{figure}[htbp]
    \centering
    % \includegraphics[width=0.8\textwidth]{runtime_plot.png} 
    \rule{0.8\textwidth}{0.3\textwidth} % This is a gray placeholder block, remove this line after uploading the image and uncomment the line above
    \caption{Average insertion time vs. Number of elements ($N$).}
    \label{fig:runtime}
\end{figure}

% ==========================================
\section{Conclusion}
% ==========================================
Summarize the findings. Reiterate how the implementation successfully maintains the balanced tree properties and achieves the theoretical $O(\log n)$ performance.

% ==========================================
% References
% ==========================================
\begin{thebibliography}{9}
\bibitem{cormen2009} 
T. H. Cormen, C. E. Leiserson, R. L. Rivest, and C. Stein, 
\textit{Introduction to Algorithms}, 3rd ed. 
Cambridge, MA, USA: MIT Press, 2009.

\bibitem{course_materials}
CSCI 4041 Course Staff, \textit{Midterm Project Module and Lectures}, University of Minnesota, Spring 2026.
\end{thebibliography}

\end{document}